%%%%% GANTT CHART 
\usepackage{pgfgantt} %<---- Gantt chart

\usepackage{pifont}
%%% IN HEADER
\usepackage{xcolor}
\definecolor{parchment}{RGB}{214, 204, 169}

\usepackage{extpfeil}
\usepackage{mdframed}
\usepackage{microtype}
\usepackage{verbatim}
\usepackage{graphicx}
\usepackage{rotating}
\usepackage{todonotes}
\usepackage{tkz-graph}
\usepackage{amsthm}
\usepackage{amsmath} % flush equations left


%for the R-N Lattice
\usepackage{svg}

\usepackage{mathtools}
%\usepackage{biblatex}
%\addbibresource{ref.bib}


\usepackage{thm-restate}


% We set up interlinked references.
\usepackage{url}
\definecolor{linkColor}{RGB}{156,78,13}
\usepackage{hyperref}
\hypersetup{
    colorlinks = true,
    allcolors = linkColor %set uniform color (def. above) for all links
    }
% We set up nicer (postscript) fonts.

\usepackage[nameinlink,capitalize]{cleveref}
\usepackage{pslatex}

% We set up paragraphs without indentation
\setlength{\parindent}{0pt}

%\setcounter{secnumdepth}{1}

% We set up version numbering of drafts
\usepackage{datetime}
\newdateformat{versiondate}{%
\THEMONTH\THEDAY}

% We set up marginal comments for drafts
\newcommand{\Ben}[1]{ \marginpar{ \begin{flushleft}\small \textbf{B:}~{#1}\end{flushleft} } }
\newcommand{\Laura}[1]{ \marginpar{ \begin{flushleft}\small \textbf{L:}~{#1}\end{flushleft} } }
\newcommand{\Will}[1]{ \marginpar{ \begin{flushleft}\small \textbf{W:}~{#1}\end{flushleft} } }


\newcommand{\openepigraph}[1]{%
  \fontsize{18}{20}\fontfamily{mathptmx}\selectfont
  \begin{fullwidth}
  \sffamily\large
  \begin{doublespace}
  \noindent\allcaps{#1}\\% epigraph
  %\noindent\allcaps{#2}% author
  \end{doublespace}
  \end{fullwidth}
}


\usepackage{algorithm}
\usepackage{algpseudocodex}
\renewcommand{\algorithmicrequire}{\textbf{Input:}}
\renewcommand{\algorithmicensure}{\textbf{Output:}}

\usepackage{verbatim}

\usepackage{framed}
\newcommand{\decisionproblem}[3]{%
\begin{framed}
\vspace{0,1cm} \noindent
% \fbox{%
% \begin{minipage}{0.96\textwidth}%
%  \begin{tabular*}{\textwidth}{@{\extracolsep{\fill}}lr} 
 \textbf{\textsc{#1} }
 % & \\ \end{tabular*}%
 %  \vspace{1.2mm}%
\par%
\noindent
{\bf{Input:}}  #2\\%
{\bf{Output:}} #3%
% \end{minipage}} \vspace{0,3cm}%
\end{framed}
}

% We set up Tikz illustration scaling and commutative diagrams
\usepackage{pgf,tikz,environ}
\usetikzlibrary{arrows}
\usetikzlibrary{cd}

% We set the bottom to ragged (fixes too much space in enumerations).
\raggedbottom

% We set up monkey-patching, used to swap theorem head from "Theorem 1.2" to "1.2. Theorem" without changing other defaults
\usepackage{xpatch} 

% We set up AMS formatting, including seamless paragraph and theorem numbering.
%\usepackage{amsfonts}
%\usepackage{amssymb}
%\usepackage{stix}

\usepackage{adjustbox}

\newcommand{\defeq}{\vcentcolon=}
\newcommand{\eqdef}{=\vcentcolon}

\setcounter{secnumdepth}{2} % Number sections and subsections only
\setcounter{tocdepth}{3}    % Include subsubsections in the ToC

%\newtheoremstyle{zoltanstyle}
%  {1em} % Space above
%  {\topsep} % Space below, usually \topsep
%  {} % Body font
%  {} % Indent amount
%  {\bfseries} % Theorem head font \bfseries
%  {.} % Punctuation after theorem head
%  {.5em} % Space after theorem head
%  {} % Theorem head spec (can be left empty, meaning `normal')
%\renewcommand{\qed}{\unskip\nobreak\quad\newline\textbf{Qed.}}

\newtheoremstyle{zoltanstyle}
  {0.2em} % Space above
  {0.2em} % Space below, usually \topsep
  {} % Body font
  {} % Indent amount
  {\bfseries} % Theorem head font \bfseries
  {.} % Punctuation after theorem head
  {.5em} % Space after theorem head
  {} % Theorem head spec (can be left empty, meaning `normal')
\renewcommand{\qed}{\hfill\textbf{Qed.}}

\newcommand{\pname}[1]{{\sc #1}}
\theoremstyle{zoltanstyle}
\swapnumbers
\xpatchcmd\swappedhead{~}{.~}{}{}
\newtheorem{body}{}
\crefname{body}{Paragraph}{Paragraphs}
\numberwithin{body}{section}
\newtheorem{claim}[body]{Claim}
\newtheorem{conjecture}[body]{Conjecture}
\newtheorem{corollary}[body]{Corollary}
\newtheorem{definition}[body]{}%{Definition}
\crefname{definition}{Paragraph}{Paragraphs}%I hate this.
\newtheorem{example}[body]{Example}
\newtheorem{addendum}[body]{Addendum}
\newtheorem{exercise}[body]{Exercise}
\newtheorem{lemma}[body]{Lemma}
\newtheorem{problem}[body]{Problem}
\newtheorem{proposition}[body]{Proposition}
\newtheorem{remark}[body]{Remark}
\newtheorem{observation}[body]{Observation}
\newtheorem{theorem}[body]{Theorem}
\newtheorem{intuition}[body]{Intuition}

\newtheorem{question}[body]{Question}

% A muted version of the brown already used for links.
\colorlet{exampleColor}{linkColor!65}

\surroundwithmdframed{example}
\surroundwithmdframed{lemma}
\surroundwithmdframed{theorem}
\surroundwithmdframed{proposition}

\surroundwithmdframed{addendum}

% We set up Tikz illustration scaling and commutative diagrams
\usepackage{pgf,tikz,environ}
\usepackage{quiver}
\usepackage{proof}
\usetikzlibrary{arrows}
\usetikzlibrary{arrows.meta,positioning,fit,calc}
\usetikzlibrary{cd}
\usetikzlibrary{positioning}
\usetikzlibrary{fadings}
\usetikzlibrary{decorations.pathreplacing}
\usetikzlibrary{decorations.pathmorphing}
\tikzset{snake it/.style={decorate, decoration=snake}}
\newcommand{\topline}
{%
  \tikz[remember picture,overlay] {%
    \draw[softGray] ([yshift=-1cm]current page.north west)
             -- ([yshift=-1cm,xshift=0.7\paperwidth]current page.north west);}
}

%~~~~~~~~~~~~~~~~~~~~~~~~~~~~~~~~~~~~~~~~~~~~~~~~~~~~~~~~~~~~~~~~~~~~~~~~~
%~~  D    theorem environs and frames  ~~~~~~~~~~~~~~~~~~~~~~~~~~~~~~~~~~~

%--  special styles for bits of text inside theorems etc  ----------------
\newcommand{\temporal}[1]{\mathsf{Temp}_{[#1]}}

\renewcommand\qedsymbol{$\blacksquare$}

\newcommand{\xmark}{\text{\ding{55}}}
\newcommand{\mycmark}{\text{\ding{51}}}

%--  shaded theorems etc  ------------------------------------------------

\definecolor{parchment}{RGB}{214, 204, 169}


\makeatletter
\renewenvironment{proof}[1][\proofname]{\par
  \pushQED{\qed}%
  \normalfont
  \topsep0pt \partopsep0pt % no space before
  \trivlist
  \item[\hskip\labelsep
        \itshape
    #1\@addpunct{.}]\ignorespaces
}{%
  \popQED\endtrivlist\@endpefalse
  \addvspace{4pt} % some space after
}
\makeatother


% We set up "Qed." proof endings.
\expandafter\let\expandafter\oldproof\csname\string\proof\endcsname
\let\oldendproof\endproof

\renewenvironment{proof}[1][\proofname]{%
  \oldproof[\normalfont \bfseries #1.]%
}{\oldendproof}

%to number and reference axioms
\usepackage[inline]{enumitem}
\usepackage{mathptmx}

\usepackage[T1]{fontenc}
\usepackage[utf8]{inputenc}

%%%%%%%%%%%%%%%%%%%%%%%%%%%%%%%%%%%%%%%%%%%%%%%%%%%%%%%%%%%%%%%%%%%%%%%%%%%%%%%%%%%%%%%%%%%%%%%%%%%%%%%%%
%%%%%%%%%%%%%%%%%%%%%%%%%%%%%%%%%%%%%%%%%%%%%%%%%%%%%%%%%%%%%%%%%%%%%%%%%%%%%%%%%%%%%%%%%%%%%%%%%%%%%%%%%
\usetikzlibrary {positioning}
\usetikzlibrary{fadings}
\usetikzlibrary{decorations.pathreplacing}
\usetikzlibrary{decorations.pathmorphing}
\tikzset{snake it/.style={decorate, decoration=snake}}
%
\newtheorem{open_problem}{Open Problem}
%for NLC commands and to have shorthand commands
\newcommand{\labdash}[2]{\stackrel{\mathclap{\normalfont\mbox{#1}}}{\text{#2}}}
\newcommand{\edge}[1]{\bullet \labdash{#1}{---} \bullet}
\newcommand{\arc}{\bullet \rightarrow \bullet}
\newcommand{\dir}[1]{\overrightarrow{#1}}
%for TikZ drawings
\definecolor{gold}{RGB}{255,215,0}
\definecolor{softBlack}{RGB}{45, 47, 49}
\definecolor{creamWhite}{RGB}{245,244,241}
\definecolor{softGray}{RGB}{220, 216, 214}
\definecolor{brick}{RGB}{232, 48, 48}

\definecolor{goodred}{HTML}{CC6677}
\definecolor{goodblue}{HTML}{332288}
\definecolor{goodyellow}{HTML}{DDCC77}
\definecolor{goodgreen}{HTML}{117733}
\definecolor{goodcyan}{HTML}{88CCEE}
\definecolor{goodwine}{HTML}{882255}
\definecolor{goodteal}{HTML}{44AA99}
\definecolor{goodolive}{HTML}{999933}
\definecolor{goodpurple}{HTML}{AA4499}

% We set up marginal comments for drafts
\newcommand{\BenSays}[1]{ \marginpar{ \begin{flushleft}\small \textbf{B:}~{#1}\end{flushleft} } }
\newcommand{\ZoltanSays}[1]{ \marginpar{ \begin{flushleft}\small \textbf{Z:}~{#1}\end{flushleft} } }

%%% ADDED BY BEN: 
\newcommand{\typesetoperator}[1]{\mathsf{#1}}

\newcommand{\unfinished}{ {\color{red} \textbf{TO DO: FINISH!}} } 

\usepackage{xspace}
%for TikZ drawings
\definecolor{gold}{RGB}{255,215,0}
\definecolor{softBlack}{RGB}{45, 47, 49}
\definecolor{creamWhite}{RGB}{245,244,241}
\definecolor{softGray}{RGB}{220, 216, 214}
\definecolor{brick}{RGB}{232, 48, 48}

\newcommand{\todowarn}{\PackageWarning{todonotes}{something to be done there.}}
\newcommand{\todocite}{{\color{red} \textbf{TO DO: CITE}\todowarn}\xspace}
\newcommand{\gentodo}[1]{\todo[backgroundcolor=blue!20!white]{#1}\todowarn}


\usepackage{enumitem}
%\usepackage{fullpage} %%%%%%%%<<<<<<<<-----------

%~~~~~~~~~~~~~~~~~~~~~~~~~~~~~~~~~~~~~~~~~~~~~~~~~~~~~~~~~~~~~~~~~~~~~~~~~
%~~  E    mathematical symbols (TYPOGRAPHY)  ~~~~~~~~~~~~~~~~~~~~~~~~~~~~~

%--  modifiers  ----------------------------------------------------------

\newcommand{\Oh}{\mathcal{O}}

\newcommand{\pr}{^\prime}
\newcommand{\prpr}{^{\prime\prime}}



%--  typesetting commands  -----------------------------------------------
\newcommand{\define}[1]{\textbf{#1}}
\newcommand{\freedag}{F_{\dag}}

\newcommand{\checkthis}[1]{{\color{red} TODO check: #1}}

\newcommand{\cat}[1]{\mathsf{#1}}
\newcommand{\typesetcategory}[1]{\cat{#1}}

\DeclareMathOperator{\sieve}{\typesetoperator{Sieve}}

\newcommand{\coprodcup}{\mathbin{\rotatebox[origin=c]{180}{$\Pi$}}}

\DeclareMathOperator{\op}{^{\typesetoperator{op}}}
\DeclareMathOperator{\tw}{\typesetoperator{tw}}
\DeclareMathOperator{\ltw}{\typesetoperator{ltw}}
\DeclareMathOperator{\domain}{\typesetoperator{dom}}
\DeclareMathOperator{\codomain}{\typesetoperator{cod}}
\DeclareMathOperator{\mono}{\typesetoperator{Mono}}
\DeclareMathOperator{\sub}{\typesetoperator{Sub}}
\DeclareMathOperator{\transred}{\typesetoperator{Tr}}
\DeclareMathOperator{\homset}{\typesetoperator{Hom}}
\DeclareMathOperator{\obj}{\typesetoperator{Ob}}
\DeclareMathOperator{\morphisms}{\typesetoperator{Mor}}
\DeclareMathOperator{\im}{\typesetoperator{im}}
\DeclareMathOperator{\imageObj}{\typesetoperator{im}}
\DeclareMathOperator{\mi}{\typesetoperator{mi}}
\DeclareMathOperator{\coim}{\typesetoperator{coim}}
\DeclareMathOperator{\coker}{\typesetoperator{coker}}
\let\ker\relax
\DeclareMathOperator{\ker}{\typesetoperator{ker}}
\DeclareMathOperator{\freecat}{\typesetoperator{Free}}
\DeclareMathOperator{\Udag}{\typesetoperator{U_\dag}}
\DeclareMathOperator{\CAT}{\typesetcategory{Cat}}
\DeclareMathOperator{\finset}{\typesetcategory{FinSet}}
\DeclareMathOperator{\finRel}{\typesetcategory{FinRel}}
\DeclareMathOperator{\set}{\typesetcategory{Set}}
\DeclareMathOperator{\Rel}{\typesetcategory{Rel}}
\DeclareMathOperator{\Matr}{\typesetcategory{Matr}}
\DeclareMathOperator{\GRhomo}{\typesetcategory{Gr}_{H}}
\DeclareMathOperator{\GRmono}{\typesetcategory{Gr}_{M}}
\DeclareMathOperator{\GRrefl}{\typesetcategory{Gr}_{Refl}}
\DeclareMathOperator{\dagcat}{\typesetcategory{DagCat}}
\DeclareMathOperator{\daggrph}{\typesetcategory{DagGr}}

%\DeclareMathOperator{\clique}{\typesetcategory{Clique}}
\DeclareMathOperator{\Sp}{\typesetcategory{Sp}}
\DeclareMathOperator{\Csp}{\typesetcategory{Csp}}
\DeclareMathOperator{\GRminor}{\typesetcategory{Gr}_{\prec}}
\DeclareMathOperator{\DGRhomo}{\typesetcategory{DGr}_{H}}
\DeclareMathOperator{\DGRmono}{\typesetcategory{DGr}_{M}}
\DeclareMathOperator{\MGRPH}{\typesetcategory{MGr}}
\DeclareMathOperator{\minors}{\typesetcategory{RGr}_{\prec}}
\DeclareMathOperator{\Rhomo}{\typesetcategory{Gr}_{RH}}
\DeclareMathOperator{\Rmono}{\typesetcategory{Gr}_{RM}}
\DeclareMathOperator{\nat}{\typesetcategory{Nat}}
\DeclareMathOperator{\HGr}{\typesetcategory{HGr}}
\DeclareMathOperator{\sd}{\mathfrak{D}}
\DeclareMathOperator{\tame}{\mathfrak{D}_{m}}
\DeclareMathOperator*{\colim}{\typesetoperator{colim}}
\DeclareMathOperator*{\lan}{\typesetcategory{Lan}}
\DeclareMathOperator*{\ran}{\typesetcategory{Ran}}
\DeclareMathOperator*{\mondiag}{\typesetcategory{monDiag}}
\DeclareMathOperator{\discnarr}{\mathbb{N}\cat{arr}}
\DeclareMathOperator*{\decomp}{\typesetcategory{Dcmp}}

\newcommand{\colFunctor}[1]{ {}_{#1}\typesetcategory{Col}}

\newcommand{\typesetproblem}[1]{\textsc{#1}}
%
\DeclareMathOperator{\p}{\typesetproblem{P}}
\DeclareMathOperator{\np}{\typesetproblem{NP}}
\DeclareMathOperator{\fpt}{\typesetproblem{FPT}}
\DeclareMathOperator{\Hcoloring}{H-\typesetproblem{Coloring}}
\DeclareMathOperator{\vc}{\typesetproblem{VertexCover}_k}
\DeclareMathOperator{\clique}{\typesetproblem{Clique}_k}
\DeclareMathOperator{\hampath}{\typesetproblem{HamiltonianDisjointPaths}}
\DeclareMathOperator{\planardeletion}{\typesetproblem{PlanarDeletion}}
\DeclareMathOperator{\chordaldeletion}{\typesetproblem{ChordalDeletion}}

\newcommand{\proarrow}{\ensuremath{\,\mathaccent\shortmid\longrightarrow\,}}



%--  doublestruck and caligraphic upper letters (first batch)  -----------

\DeclareMathOperator{\Aa}{\mathcal{A}}  \DeclareMathOperator{\Ab}{\mathbb{A}}
\DeclareMathOperator{\Ba}{\mathcal{B}}  \DeclareMathOperator{\Bb}{\mathbb{B}}
\DeclareMathOperator{\Ca}{\mathcal{C}}  \DeclareMathOperator{\Cb}{\mathbb{C}}
\DeclareMathOperator{\Da}{\mathcal{D}}  \DeclareMathOperator{\Db}{\mathbb{D}}
\DeclareMathOperator{\Ea}{\mathcal{E}}  \DeclareMathOperator{\Eb}{\mathbb{E}}
\DeclareMathOperator{\Fa}{\mathcal{F}}  \DeclareMathOperator{\Fb}{\mathbb{F}}
\DeclareMathOperator{\Ga}{\mathcal{G}}  \DeclareMathOperator{\Gb}{\mathbb{G}}
\DeclareMathOperator{\Ha}{\mathcal{H}}  \DeclareMathOperator{\Hb}{\mathbb{H}}
\DeclareMathOperator{\Ia}{\mathcal{I}}  
\DeclareMathOperator{\Ja}{\mathcal{J}}  
\DeclareMathOperator{\Ka}{\mathcal{K}}  \DeclareMathOperator{\Kb}{\mathbb{K}}
\DeclareMathOperator{\La}{\mathcal{L}}  \DeclareMathOperator{\Lb}{\mathbb{L}}
\DeclareMathOperator{\Ma}{\mathcal{M}}  \DeclareMathOperator{\Mb}{\mathbb{M}}
\DeclareMathOperator{\Na}{\mathcal{N}}  \DeclareMathOperator{\Nb}{\mathbb{N}}
\DeclareMathOperator{\Oa}{\mathcal{O}}  \DeclareMathOperator{\Ob}{\mathbb{O}}
\DeclareMathOperator{\Pa}{\mathcal{P}}  \DeclareMathOperator{\Pb}{\mathbb{P}}
\DeclareMathOperator{\Qa}{\mathcal{Q}}  \DeclareMathOperator{\Qb}{\mathbb{Q}}
\DeclareMathOperator{\Ra}{\mathcal{R}}  \DeclareMathOperator{\Rb}{\mathbb{R}}
\DeclareMathOperator{\Sa}{\mathcal{S}}  \DeclareMathOperator{\Sb}{\mathbb{S}}
\DeclareMathOperator{\Ta}{\mathcal{T}}  \DeclareMathOperator{\Tb}{\mathbb{T}}
\DeclareMathOperator{\Ua}{\mathcal{U}}  \DeclareMathOperator{\Ub}{\mathbb{U}}
\DeclareMathOperator{\Va}{\mathcal{V}}  \DeclareMathOperator{\Vb}{\mathbb{V}}
\DeclareMathOperator{\Wa}{\mathcal{W}}  \DeclareMathOperator{\Wb}{\mathbb{W}}
\DeclareMathOperator{\Xa}{\mathcal{X}}  \DeclareMathOperator{\Xb}{\mathbb{X}}
\DeclareMathOperator{\Ya}{\mathcal{Y}}  \DeclareMathOperator{\Yb}{\mathbb{Y}}
\DeclareMathOperator{\Za}{\mathcal{Z}}  \DeclareMathOperator{\Zb}{\mathbb{Z}}

%--  doublestruck and caligraphic upper letters (second batch)  ----------

%\newcommand{\Aa}{\mathbb{A}}\newcommand{\aA}{\mathcal{A}}
%\newcommand{\Bb}{\mathbb{B}}\newcommand{\bB}{\mathcal{B}}
\newcommand{\Cc}{\mathbb{C}}\newcommand{\cC}{\mathcal{C}}
\newcommand{\Dd}{\mathbb{D}}\newcommand{\dD}{\mathcal{D}}
\newcommand{\Ee}{\mathbb{E}}\newcommand{\eE}{\mathcal{E}}
\newcommand{\Ff}{\mathbb{F}}\newcommand{\fF}{\mathcal{F}}
\newcommand{\Gg}{\mathbb{G}}\newcommand{\gG}{\mathcal{G}}
\newcommand{\Hh}{\mathbb{H}}\newcommand{\hH}{\mathcal{H}}
\newcommand{\Ii}{\mathbb{I}}\newcommand{\iI}{\mathcal{I}}
\newcommand{\Jj}{\mathbb{J}}\newcommand{\jJ}{\mathcal{J}}
\newcommand{\Kk}{\mathbb{K}}\newcommand{\kK}{\mathcal{K}}
\newcommand{\Ll}{\mathbb{L}}\newcommand{\lL}{\mathcal{L}}
\newcommand{\Mm}{\mathbb{M}}\newcommand{\mM}{\mathcal{M}}
\newcommand{\Nn}{\mathbb{N}}\newcommand{\nN}{\mathcal{N}}
\newcommand{\Oo}{\mathbb{O}}\newcommand{\oO}{\mathcal{O}}
\newcommand{\Pp}{\mathbb{P}}\newcommand{\pP}{\mathcal{P}}
\newcommand{\Qq}{\mathbb{Q}}\newcommand{\qQ}{\mathcal{Q}}
\newcommand{\Rr}{\mathbb{R}}\newcommand{\rR}{\mathcal{R}}
\newcommand{\Ss}{\mathbb{S}}\newcommand{\sS}{\mathcal{S}}
\newcommand{\Tt}{\mathbb{T}}\newcommand{\tT}{\mathcal{T}}
\newcommand{\Uu}{\mathbb{U}}\newcommand{\uU}{\mathcal{U}}
\newcommand{\Vv}{\mathbb{V}}\newcommand{\vV}{\mathcal{V}}
\newcommand{\Ww}{\mathbb{W}}\newcommand{\wW}{\mathcal{W}}
\newcommand{\Xx}{\mathbb{X}}\newcommand{\xX}{\mathcal{X}}
\newcommand{\Yy}{\mathbb{Y}}\newcommand{\yY}{\mathcal{Y}}
\newcommand{\Zz}{\mathbb{Z}}\newcommand{\zZ}{\mathcal{Z}}

%--  sans serif and frak lower letters  ----------------------------------

\newcommand{\sfa}{\mathsf{a}}\newcommand{\fra}{\mathcal{a}}
\newcommand{\sfb}{\mathsf{b}}\newcommand{\frb}{\mathcal{b}}
\newcommand{\sfc}{\mathsf{c}}\newcommand{\frc}{\mathcal{c}}
\newcommand{\sfd}{\mathsf{d}}\newcommand{\frd}{\mathcal{d}}
\newcommand{\sfe}{\mathsf{e}}\newcommand{\fre}{\mathcal{e}}
\newcommand{\sff}{\mathsf{f}}\newcommand{\frf}{\mathcal{f}}
\newcommand{\sfg}{\mathsf{g}}\newcommand{\frg}{\mathcal{g}}
\newcommand{\sfh}{\mathsf{h}}\newcommand{\frh}{\mathcal{h}}
\newcommand{\sfi}{\mathsf{i}}\newcommand{\fri}{\mathcal{i}}
\newcommand{\sfj}{\mathsf{j}}\newcommand{\frj}{\mathcal{j}}
\newcommand{\sfk}{\mathsf{k}}\newcommand{\frk}{\mathcal{k}}
\newcommand{\sfl}{\mathsf{l}}\newcommand{\frl}{\mathcal{l}}
\newcommand{\sfm}{\mathsf{m}}\newcommand{\frm}{\mathcal{m}}
\newcommand{\sfn}{\mathsf{n}}\newcommand{\frn}{\mathcal{n}}
\newcommand{\sfo}{\mathsf{o}}\newcommand{\fro}{\mathcal{o}}
\newcommand{\sfp}{\mathsf{p}}\newcommand{\frp}{\mathcal{p}}
\newcommand{\sfq}{\mathsf{q}}\newcommand{\frq}{\mathcal{q}}
\newcommand{\sfr}{\mathsf{r}}\newcommand{\frr}{\mathcal{r}}
\newcommand{\sfs}{\mathsf{s}}\newcommand{\frs}{\mathcal{s}}
\newcommand{\sft}{\mathsf{t}}\newcommand{\frt}{\mathcal{t}}
\newcommand{\sfu}{\mathsf{u}}\newcommand{\fru}{\mathcal{u}}
\newcommand{\sfv}{\mathsf{v}}\newcommand{\frv}{\mathcal{v}}
\newcommand{\sfw}{\mathsf{w}}\newcommand{\frw}{\mathcal{w}}
\newcommand{\sfx}{\mathsf{x}}\newcommand{\frx}{\mathcal{x}}
\newcommand{\sfy}{\mathsf{y}}\newcommand{\fry}{\mathcal{y}}
\newcommand{\sfz}{\mathsf{z}}\newcommand{\frz}{\mathcal{z}}


%%% Lineale stuff %%%

\newcommand{\product}{\&}
\newcommand{\coproduct}{\oplus}
\newcommand{\productmap}[2]{\langle #1, #2\rangle}
\newcommand{\coproductmap}[2]{[#1,#2]}
\newcommand{\eval}{\mathsf{eval}}
\newcommand{\discard}{\epsilon} % discarding map in Set
\newcommand{\tensor}{\otimes}
\newcommand{\unit}{I}
\newcommand{\inthom}[2]{[#1,#2]}
\newcommand{\linorder}{\sqsubseteq}
\newcommand{\lintensor}{\otimes}
\newcommand{\lininthom}{\multimap} % internal hom


%%% Time intervals %%%

\newcommand{\interval}[2]{#1 \downarrow P \downarrow  #2}

%%Added by Fred:

\newcommand{\grp}{\cat{Grph}}
\newcommand{\cu}[2]{\cat{Cu}\left(#1, #2\right)}
\newcommand{\tpfrak}[1]{\mathfrak{P} \circ #1}
\newcommand{\IN}{\cat{I_{\Nb}}}
\newcommand{\rea}{\rightsquigarrow}

\newcommand{\Cov}{{\mathsf{Cov}}}
\newcommand{\sh}{\typesetoperator{sh}}


% Laura:
\newcommand{\tcdtext}{tree co-decomposition\xspace}
\newcommand{\mfdiag}{matching family diagram\xspace}
%----------
\newcommand{\tcal}{\ensuremath{\mathcal{T}}\xspace}
\newcommand{\tcdfull}{\ensuremath{(T,\{G_a\},\{g^e_t\})}}

%universal extension
\newcommand{\uext}[1]{\mathbf{ext}{#1}}
%local extension
\newcommand{\ext}[2]{\mathbf{ext}_{#1}{#2}}

\DeclareMathOperator{\children}{\mathbf{children}}
%\DeclareMathOperator*{\lop}{\Diamond}
\newcommand{\lop}[1]{%
  \mathop{\scalebox{1.5}{$\Diamond$}_{\!#1}}\displaylimits
}

\DeclareMathOperator{\lc}{\mathbf{left}}
\DeclareMathOperator{\rc}{\mathbf{right}}

\DeclareMathOperator{\atoms}{\mathbb{A}}

\newcommand{\bigotimesalpha}{%
  \mathop{\sideset{}{_{\alpha}}\bigotimes}\displaylimits
}

\newcommand{\bigO}[1]{{\mathcal{O}\!\left(#1\right)}}

\newcommand{\pts}[1]{\underline{#1}}